\section{The exact cycle index and support constants}\label{sec:identity}
For a partition $\nu$ with $m_\ell(\nu)$ parts equal to $\ell$, define
\[
z_\nu=\prod_{\ell\geq1}\ell^{m_\ell(\nu)}m_\ell(\nu)!,\qquad
p_\nu=\prod_{\ell\geq1}p_\ell^{m_\ell(\nu)},\qquad
h_j(p)=\sum_{\nu\vdash j}\frac{p_\nu}{z_\nu}.
\]
Let $r(\nu)$ be the number of square roots in $S_{|\nu|}$ of a fixed permutation of cycle type $\nu$. This is a square-root count, not a count of involutions commuting with that permutation. With
\begin{equation}\label{eq:F}
F(p)=\prod_{j\geq1}(1-h_j(p))^{-1},\qquad
\alpha_\nu=[p_\nu]F(p),
\end{equation}
Schwob's identity is
\begin{equation}\label{eq:schwob}
a_n=\sum_{\nu\vdash n}\alpha_\nu r(\nu)
 =\sum_{\mu\vdash n}\frac1{|S_\mu|}\sum_{\sigma\in S_\mu}r(\sigma).
\end{equation}
Here $S_\mu$ is the standard Young subgroup, a product of symmetric groups on the blocks specified by $\mu$. All coefficients of $F$ are nonnegative. Formula~\eqref{eq:schwob} is imported from \cite[\S4.3]{Schwob}; it fixes the normalization of the entire argument.

\subsection{Square roots and nonfixed support}
Write $\nu=1^{n-s}\rho$, where $\rho$ has no parts equal to $1$ and $|\rho|=s$. Then
\begin{equation}\label{eq:rootfactor}
r(1^{n-s}\rho)=I_{n-s}r(\rho).
\end{equation}
Indeed, a square root commutes with its square and therefore preserves the fixed-point set of its square. On that set it is an arbitrary involution; on the complement it is a square root of the permutation of type $\rho$.

For later estimates and computations, the contribution of $m$ cycles of length $\ell$ to $r(\nu)$ is
\begin{equation}\label{eq:rootcycles}
\begin{cases}
\displaystyle\sum_{k=0}^{\lfloor m/2\rfloor}
 \frac{m!\ell^k}{(m-2k)!2^kk!},&\ell\text{ odd},\\[7pt]
0,&\ell\text{ even and }m\text{ odd},\\[3pt]
\displaystyle\frac{m!\ell^{m/2}}{2^{m/2}(m/2)!},
 &\ell\text{ even and }m\text{ even}.
\end{cases}
\end{equation}
The factors are multiplied over $\ell$. One may pair two cycles of the same length, with $\ell$ possible interlacings; an unpaired odd cycle has a unique root of the same length, while an even cycle cannot be left unpaired. This also explains the source formula in \cite[\S3.1]{Schwob}.

Let $u_\ell$ have weight $\ell$ for $\ell\geq2$, and write $u_\rho=\prod_{\ell\geq2}u_\ell^{m_\ell(\rho)}$. For $j\geq2$ set
\begin{equation}\label{eq:dj}
d_j(u)=\sum_{\substack{\rho\text{ has parts }\geq2\\0<|\rho|\leq j}}
 \frac{\fall{j}{|\rho|}}{(j!-1)z_\rho}\,u_\rho,
\qquad \fall{j}{s}=\frac{j!}{(j-s)!}.
\end{equation}
Define the convergent fixed-weight coefficient product
\begin{equation}\label{eq:P}
P(u)=\prod_{j\geq2}(1-d_j(u))^{-1},\qquad
q_\rho=[u_\rho]P(u),
\end{equation}
and put
\begin{equation}\label{eq:Hs}
H_s=\sum_{\substack{\rho\vdash s\\\text{parts }\geq2}}r(\rho)q_\rho,
\qquad H_0=1.
\end{equation}
The empty coefficient of $P$ is $1$. The sum for $H_1$ is empty, and the only type at support $2$ is $(2)$, which has no square root. Hence $H_1=H_2=0$.

The notation $q_\rho$ will also be useful for the companions. In particular, $q_{(2)}$ is positive, although $H_2=0$. Keeping these two quantities separate avoids confusing a cycle-index coefficient with its root-weighted sum. Absolute convergence, finite computation, and coefficient stabilization are proved next.
